\documentclass[10pt,a4paper]{amsart}
\usepackage[margin=19mm]{geometry}
\usepackage{amsmath,amssymb,mathtools}
\usepackage[scr=rsfs,cal=euler]{mathalfa}
\usepackage{xfrac}
\usepackage[unicode,hidelinks,bookmarks=false]{hyperref}
\newcommand{\ie}{i.e.\ }
\newcommand{\cf}{cf.\ }
\newcommand{\resp}{respectively\ }
\newcommand{\Subsection}{\subsection}
\providecommand{\nobreakdash}{\nobreak\hbox{-}}
\setlength{\emergencystretch}{2em}
\multlinegap0pt
\usepackage{dsfont}
\usepackage{amssymb}
\usepackage{float}
\usepackage{bm}
\usepackage{tikz}
\newtheorem{definition}{Definition}
\newtheorem{theorem}{Theorem}
\title{Equilibria for Time-inconsistent Regular-singular Control Problems}
\author{Yuting Jia, Zongxia Liang, Xiaodong Luo}
\date{Source: arXiv:2609.08877v1 (CC BY 4.0)}
\begin{document}
\maketitle

\noindent\emph{Source and licence.} Equilibria for Time-inconsistent Regular-singular Control Problems. Version arXiv:2609.08877v1 (CC BY 4.0), distributed by arXiv under the Creative Commons Attribution 4.0 International licence, http://creativecommons.org/licenses/by/4.0/, as recorded in the arXiv licence metadata for this exact version and shown on the abstract page https://arxiv.org/abs/2609.08877v1. Full licence notice: \texttt{licenses/arxiv-2609.08877-CC-BY-4.0.txt}.\par\smallskip

\noindent\textit{Editorial scope.} Counted unit: the article's theorem verificationthm (source0.tex lines 271--340). The article's complete original proof of it is delivered with it (source0.tex lines 342--473, one \texttt{proof} environment of 132 lines). No mathematics is written, repaired or re-derived: the statement and the proof are the article's own lines, in the article's own notation, and no retained line is substituted or re-flowed.  Retained uncounted as the source-local context the proof needs: lines 130--137; lines 146--146; lines 180--220; lines 249--252.  Nothing was omitted from any retained span, so the package carries no omission record.  No retained line contains a citation, so the wrapper carries no bibliography and no bibliography entry.  No retained line contains a figure environment, an image-inclusion command or drawing code, so no image is delivered and the package declares no asset dependency.  Editorial formatting only, in addition: an AMS \texttt{amsart} preamble that restates the article's own declarations and macros used by the retained lines, namely the environments \texttt{definition}, \texttt{theorem} on separate counters, together with the standard packages those lines need.  The retained environments print in this document as Definition 1, Theorem 1 in order of appearance, whereas the article prints the same items with the numbering of its own full document.  The complete native source of the article as distributed in the arXiv e-print for this exact version is bundled unchanged as \texttt{sources/209759/source0.tex}.
For any fixed $(x,t,y)\in \mathcal{Q}:=[0,+\infty)\times[0,+\infty)\times[0,+\infty)$, the controlled state process $X^{\pi,\xi}$ under the regular control $\pi$ and the singular control $\xi$ follows the general stochastic differential equation (abbr. SDE):
\begin{equation}
\left\{\begin{array}{l}
\mathrm{d}X^{\pi,\xi}_{r}=\mu\left(X^{\pi,\xi}_{r},r,\pi_{r},\xi_{r}\right)\mathrm{d}r+\sigma\left(X^{\pi,\xi}_{r},r,\pi_{r},\xi_{r}\right)\mathrm{d}B_{r}-\mathrm{d}\xi_{r},\ r\in[t,\tau^{x,t,y;\pi,\xi}_{t}]\footnote{In this paper, for any stopping time $\tau$, when $\tau=+\infty$, $[a,\tau]$ represents $[a,+\infty)$ for any $a\in\mathbb{R}$.},\\
X^{\pi,\xi}_{t-}=x,
\end{array}\right.\label{sde-regular}
\end{equation}
where $\pi=\left\{\pi_r\right\}_{r\ge t}$ is the regular control, $\xi=\left\{\xi_r\right\}_{r\ge t}$ is the singular control which is nondecreasing and c$\acute{a}$dl$\acute{a}$g with initial $\xi_{t-}=y$, and $\Delta \xi_r=\xi_r-\xi_{r-}\le X_{r-}^{\pi,\xi}$. $\tau^{x,t,y;\pi,\xi}_{t}:=\inf\left\{r\ge t|X^{\pi,\xi}_{r}= 0\right\}$ is the time of bankruptcy, and the drift function $\mu$  and the volatility function $\sigma$ are exogenously given and deterministic continuous functions on $[0,+\infty)\times[0,+\infty)\times U\times[0,+\infty)$, where compact set $U\subset\mathbb{R}$ is the regular control space.%U Borel-measurable, U compact 为了验证定理过得去，例子也是满足，不过也可以加上扰动过程u右连续（没必要右连左极），但是我们的例子的均衡策略不是右连续的，而且不想加扰动过程u右连续的条件。也可以对于H关于\pi变量的性质加条件。

\begin{equation}c(x-a,t,y+a)=c(x,t,y),\ \forall (x,t,y)\in \mathcal{Q}, a\geq 0,\label{cxy}\end{equation} or, equivalently, $c_x(x,t,y)=c_y(x,t,y), \forall (x,t,y)\in \mathcal{Q}$.

\begin{definition}[Admissible regular-singular control pair]\label{admissible2}
Fix initial $(x,t,y)\in\mathcal{Q}$. Let $\pi=\{\pi_{r}\}_{r\ge t}$ be a $U$-valued $\{\mathcal{F}_{r}\}_{r\ge t}$-progressively measurable process and $\xi=\{\xi_{r}\}_{r\ge t}$ be a nondecreasing c$\acute{a}$dl$\acute{a}$g $\{\mathcal{F}_{r}\}_{r\ge t}$-adapted process. Then $(\pi,\xi)$ is called an admissible regular-singular control pair at $(x,t,y)$ if the following hold:\\
(a). $\xi_{t-}=y$, $\Delta \xi_r=\xi_r-\xi_{r-}\le X_{r-}^{\pi,\xi}$, and the SDE \eqref{sde-regular} has a unique strong solution $X^{\pi,\xi}=\{X^{\pi,\xi}_{r}\}_{r\in [t,\tau^{x,t,y;\pi,\xi}_{t}]}$.\\
(b). The strong solution $X^{\pi,\xi}$ in (a) and $(\pi,\xi)$ satisfies
\begin{equation}
\begin{aligned}
&\mathbb{E}_{x, t, y}\left[\int_t^{\tau^{x,t,y;\pi,\xi}_{t}} \beta(r-t) \left|H\left(X_r^{\pi,\xi}, r, \pi_r, \xi_r \right)\right| \mathrm{d} r\right]<+\infty,\\
&\mathbb{E}_{x, t, y}\left[\int_t^{\tau^{x,t,y;\pi,\xi}_{t}} \beta(r-t) \left|c\left(X_{r-}^{\pi,\xi}, r, \xi_{r-}\right)\right| \mathrm{d} \xi_r\right]<+\infty.
\end{aligned}\nonumber
\end{equation}
Denote the set of all admissible regular-singular control pairs at $(x,t,y)$ by $\mathcal{D}^{x,t,y}$. 

For any fixed admissible regular-singular control law $\left(\hat\Pi, \hat\Xi\right)$, for any $(\eta,u)\in\mathcal{D}^{x,t,y}$, if there exist constants $\tilde{h}>0$ and $\tilde{M}>0$ such that $\eta$ additionally satisfies
\begin{equation}
\label{growthcond}
\eta_{(t+h)-}-\eta_{t}\le \tilde{M}h,\ \text{a.s.},\ \forall h\in(0,\tilde{h}),
\end{equation}
and the perturbed regular-singular control pair $\left(\pi^{h},\xi^{h}\right)$, $\forall h\in(0,\tilde{h})$, defined by
\begin{equation}
\begin{aligned}
&\pi^{h}_{r}:=\begin{cases}
u_{r}, & r\in[t,\tau^{x,t,y;u,\eta}_{t}\wedge((t+h)-)]\footnote{In this paper, when $\tau^{x,t,y;u,\eta}_{t}<t+h$, $[t,\tau^{x,t,y;u,\eta}_{t}\wedge((t+h)-)]$ refers to $[t,\tau^{x,t,y;u,\eta}_{t}]$, and when $\tau^{x,t,y;u,\eta}_{t}\ge t+h$, $[t,\tau^{x,t,y;u,\eta}_{t}\wedge((t+h)-)]$ refers to $[t,t+h)$.},\\
\pi^{X^{u,\eta}_{\tau^{x,t,y;u,\eta}_{t}\wedge((t+h)-)},\tau^{x,t,y;u,\eta}_{t}\wedge(t+h),\eta_{\tau^{x,t,y;u,\eta}_{t}\wedge((t+h)-)};\hat{\Pi},\hat{\Xi}}_{r}, & r\in(\tau^{x,t,y;u,\eta}_{t}\wedge((t+h)-),\tau^{x,t,y;\pi^h,\xi^h}_{t}]\footnote{Similarly, when $\tau^{x,t,y;u,\eta}_{t}<t+h$, $(\tau^{x,t,y;u,\eta}_{t}\wedge((t+h)-),\tau^{x,t,y;\pi^h,\xi^h}_{t}]$ refers to $(\tau^{x,t,y;u,\eta}_{t},\tau^{x,t,y;\pi^h,\xi^h}_{t}]=(\tau^{x,t,y;u,\eta}_{t},\tau^{x,t,y;u,\eta}_{t}]=\emptyset$, and when $\tau^{x,t,y;u,\eta}_{t}\ge t+h$, $(\tau^{x,t,y;u,\eta}_{t}\wedge((t+h)-),\tau^{x,t,y;\pi^h,\xi^h}_{t}]$ refers to $[t+h,\tau^{x,t,y;\pi^h,\xi^h}_{t}]$.},
\end{cases}\\
&\xi^{h}_{r}:=\begin{cases}
\eta_{r}, & r\in[t,\tau^{x,t,y;u,\eta}_{t}\wedge((t+h)-)],\\
\xi^{X^{u,\eta}_{\tau^{x,t,y;u,\eta}_{t}\wedge((t+h)-)},\tau^{x,t,y;u,\eta}_{t}\wedge(t+h),\eta_{\tau^{x,t,y;u,\eta}_{t}\wedge((t+h)-)};\hat{\Pi},\hat{\Xi}}_{r}, &  r\in(\tau^{x,t,y;u,\eta}_{t}\wedge((t+h)-),\tau^{x,t,y;\pi^h,\xi^h}_{t}],
\end{cases}
\end{aligned}\label{pihxih}
\end{equation}
and the corresponding state process
\begin{equation}
\begin{aligned}
X^{\pi^{h},\xi^{h}}_{r}=\begin{cases}
X^{u,\eta}_{r}, & r\in[t,\tau^{x,t,y;u,\eta}_{t}\wedge((t+h)-)],\\
X^{X^{u,\eta}_{\tau^{x,t,y;u,\eta}_{t}\wedge((t+h)-)},\tau^{x,t,y;u,\eta}_{t}\wedge(t+h),\eta_{\tau^{x,t,y;u,\eta}_{t}\wedge((t+h)-)};\hat{\Pi},\hat{\Xi}}_{r}, &  r\in(\tau^{x,t,y;u,\eta}_{t}\wedge((t+h)-),\tau^{x,t,y;\pi^h,\xi^h}_{t}],
\end{cases}
\end{aligned}\label{xpihxih}
\end{equation}
satisfy integrability condition (b), where $\tau^{x,t,y;\pi^h,\xi^h}_{t}:=\inf\left\{r\ge t|X^{\pi^h,\xi^h}_{r}= 0\right\}$, then we call $\left(\eta,u\right)$ an admissible pair of perturbations at $(x,t,y)$  w.r.t. $\left(\hat\Pi, \hat\Xi\right)$. We denote the set of all admissible pairs of perturbations at $(x,t,y)$  w.r.t. $\left(\hat\Pi, \hat\Xi\right)$ by $\mathcal{\tilde D}^{x,t,y;\hat\Pi,\hat\Xi}$.
\end{definition}

\begin{equation}
\label{equicond}
\liminf\limits_{h\downarrow 0}\frac{J\left(x,t,y;\pi^{h},\xi^{h}\right)-J\left(x,t,y;\hat{\pi},\hat{\xi}\right)}{h}\ge 0,
\end{equation}

\begin{theorem}[Verification Theorem]\label{verificationthm}
Given a function $V(x,t,y)$ and a family of functions $\left\{f^{s}(x,t,y)\right\}_{s\in[0,t]}$. Define $f(x,t,y,s):=f^{s}(x,t,y)$. Define the infinitesimal operator
\begin{equation}
\begin{aligned}
&\left(\mathcal{A}^{u}\varphi\right)(x,t,y):=\varphi_{t}(x,t,y)+\mu(x,t,u,y)\varphi_{x}(x,t,y)+\frac{1}{2}\sigma^{2}(x,t,u,y)\varphi_{xx}(x,t,y),\ \forall\varphi:\mathcal{Q}\rightarrow\mathbb{R},\\
&\left(\mathcal{A}^{\Pi}\varphi\right)(x,t,y):=\left(\mathcal{A}^{\Pi(x,t,y)}\varphi\right)(x,t,y),\ \forall\varphi:\mathcal{Q}\rightarrow\mathbb{R}.
\end{aligned}\nonumber
\end{equation} 
Assume the following conditions hold:\\
(1).
\begin{equation}
\begin{aligned}
& V \in C^{2,1,1}\left(\mathcal{Q}\right), \\
& f \in C^{2,1,1,1}\left(\left\{(x,t,y,s)\mid (x,t,y)\in\mathcal{Q}, s\in [0,t]\right\}\right).
\end{aligned}\nonumber
\end{equation}
(2). There exists an admissible regular-singular control law $\left(\hat{\Pi},\hat{\Xi}\right)$ such that the singular control law $\hat{\Xi}$ satisfies
\begin{equation}
\begin{aligned}
&W^{\hat{\Xi}}:=\left\{(x,t,y)\in\mathcal{Q}|c(x,t,y)-V_{x}(x,t,y)+V_{y}(x,t,y)>0\right\},\\
&P^{\hat{\Xi}}:=\left\{(x,t,y)\in\mathcal{Q}|c(x,t,y)-V_{x}(x,t,y)+V_{y}(x,t,y)=0\right\},
\end{aligned}\nonumber
\end{equation}
and the regular control law $\hat{\Pi}$ satisfies
\begin{equation}
\hat{\Pi}(x,t,y)\in\mathop{\mathrm{argmin}}\limits_{u\in U}\left\{\left(\mathcal{A}^{u}V\right)(x,t,y)+H(x,t,u,y)-f_s(x,t,y,t)\right\},\ \forall (x,t,y)\in\mathcal{Q},
\nonumber
\end{equation}
$\hat{\Pi}$ is continuous on $\overline{W^{\hat{\Xi}}}-W^{\hat{\Xi}}$ when $\hat{\Pi}$ is regarded as a function on $\overline{W^{\hat{\Xi}}}$. \\
(3). $V(x,t,y)$ and $f^{s}(x,t,y)$ satisfy

\begin{subequations}
\begin{align}
&\min\left\{\left(\mathcal{A}^{\hat{\Pi}}V\right)(x,t,y)+H\left(x,t,\hat{\Pi}(x,t,y),y\right)-f_s(x,t,y,t),\right.\nonumber\\
&\left.\quad \quad c(x,t,y)-V_{x}(x,t,y)+V_{y}(x,t,y)\right\}=0,\ \forall (x,t,y)\in\mathcal{Q}, \label{v}\\
&V(0,t,y)=0,\ \forall(0,t,y)\in\mathcal{Q},\label{vT}\\
&\left(\mathcal{A}^{\hat{\Pi}}f^{s}\right)(x,t,y)+\beta(t-s)H\left(x,t,\hat\Pi(x,t,y),y\right)=0,\ \forall (x,t,y)\in \overline{W^{\hat{\Xi}}}, s\in[0,t],\label{fw}\\
& \beta(t-s)c(x,t,y)-f^{s}_{x}(x,t,y)+f^{s}_{y}(x,t,y)=0,\ \forall (x,t,y)\in P^{\hat{\Xi}}, s\in[0,t],\label{fp}\\
&f^{s}(0,t,y)=0,\ \forall (0,t,y)\in\mathcal{Q}, s\in[0,t].\label{fT}
\end{align}
\end{subequations}
(4). For functions $\varphi=f^{s},V,(x,t,y)\mapsto f(x,t,y,t)$, for any $(x,t,y)\in\mathcal{Q}$, any  $s\in[0,t]$, and any $n\geq t$, it holds that 
\begin{equation}
\mathbb{E}_{x,t,y}\int_{t}^{\tau_n}\left[\varphi_{x}\left(\hat{X}_{r},r,\hat\xi_{r})\sigma(\hat{X}_{r},r,\hat\pi_r,\hat\xi_{r}\right)\right]^{2}\mathrm{d}r<+\infty,
\label{martingale}\end{equation}
where $\left(\hat{\pi},\hat{\xi}\right):=\left(\pi^{x,t,y;\hat{\Pi},\hat{\Xi}},\xi^{x,t,y;\hat{\Pi},\hat{\Xi}}\right)$ is the regular-singular control process generated by $\left(\hat{\Pi},\hat{\Xi}\right)$ at $(x,t,y)$, $\hat{X}:=X^{x,t,y;\hat{\Pi},\hat{\Xi}}$ is the corresponding state process, and $\tau_{n}:=\tau^{x,t,y;\hat{\Pi},\hat{\Xi}}_{t}\wedge n$.\\
(5). For any $(x,t,y)\in\mathcal{Q}$, any $s\in[0,t]$, it holds that 

\begin{align}
&\lim\limits_{n\rightarrow+\infty}\mathbb{E}_{x,t,y}f^{s}\left(\hat{X}_{\tau_{n}},\tau_{n},\hat{\xi}_{\tau_{n}}\right)=0,\label{finfty1}\\
&\lim\limits_{n\rightarrow+\infty}\mathbb{E}_{x,t,y}\left[V\left(\hat{X}_{\tau_{n}},\tau_{n},\hat{\xi}_{\tau_{n}}\right)-f\left(\hat{X}_{\tau_{n}},\tau_{n},\hat{\xi}_{\tau_{n}},\tau_{n}\right)\right]=0.\label{finfty2}
\end{align}
(6). For any $(x,t,y)\in\mathcal{Q}$, any admissible pair of perturbations $(u,\eta)\in\mathcal{\tilde D}^{x,t,y;\hat\Pi,\hat\Xi}$, there exists a constant $h_0>0$ such that

\begin{align}
&\mathbb{E}_{x,t,y}\int_{t}^{\tau^{x,t,y;u,\eta}_{t}\wedge(t+h_0)}\left[V_{x}\left(X^{u,\eta}_{r},r,\eta_{r})\sigma(X^{u,\eta}_{r},r,u_r,\eta_{r}\right)\right]^{2}\mathrm{d}r<+\infty,\label{martingale2}\\
&\mathbb{E}_{x,t,y}\left[\sup\limits_{\substack{r\in(t,\tau^{x,t,y;u,\eta}_{t}\wedge(t+h_{0}))\\ h\in(0,h_{0})}}\left|f_s\left(X^{u,\eta}_{r},r,\eta_{r},r\right)+(\beta(r-t)-1) H\left(X^{u,\eta}_r, r,u_r, \eta_r \right)\right.\right.\nonumber\\
&\quad\quad\left.\left.-f_s\left(X^{u,\eta}_{\tau^{x,t,y;u,\eta}_{t}\wedge((t+h)-)},\tau^{x,t,y;u,\eta}_{t}\wedge(t+h),\eta_{\tau^{x,t,y;u,\eta}_{t}\wedge((t+h)-)},r\right)\right|\right]<+\infty,\label{integrability}\\
&\mathbb{E}_{x,t,y}\left[\sup\limits_{r\in(t,\tau^{x,t,y;u,\eta}_{t}\wedge(t+h_{0}))}\left|(\beta(r-t)-1)c\left(X^{u,\eta}_{r-},r,\eta_{r-}\right)\right|\right]<+\infty.\label{integrability2}
\end{align}%\eqref{integrability2}中a\in[0,\Delta\eta_r]因为\eqref{cxy}可以去掉,之前的条件中a\in[0,\Delta\eta_r]不能去掉。

Then $\left(\hat{\Pi},\hat{\Xi}\right)$ is an equilibrium regular-singular control law and $V$ is the corresponding equilibrium value function. Moreover, $f$ has the probabilistic interpretation
\begin{equation}
\begin{aligned}
\label{intf}
f(x,t,y,s)=&\mathbb{E}_{x,t,y}\left[\int_t^{\tau^{x,t,y;\hat{\Pi},\hat{\Xi}}_{t}} \beta(r-s) H\left(\hat{X}_r, r, \hat\pi_r, \hat{\xi}_r \right) \mathrm{d} r\right. \\
& \left.+\int_t^{\tau^{x,t,y;\hat{\Pi},\hat{\Xi}}_{t}} \beta(r-s) c\left(\hat{X}_{r-}, r, \hat{\xi}_{r-}\right) \mathrm{d} \hat{\xi}_r \right],\ \forall (x,t,y)\in\mathcal{Q}, s\in[0,t].
\end{aligned}
\end{equation}
\end{theorem}

\begin{proof}
We prove Theorem \ref{verificationthm} in three steps.

{\bf Step 1:} We show that $f$ has the probabilistic interpretation \eqref{intf}.

Fix any $(x,t,y)\in\mathcal{Q}$, $s\in[0,t]$, and $n\ge t$, applying the It\^{o}--Tanaka--Meyer formula to $f^{s}$, we obtain
\begin{equation}
\begin{aligned}
&f^{s}\left(\hat{X}_{\tau_{n}},\tau_{n},\hat{\xi}_{\tau_{n}}\right)-f^{s}(x,t,y)\\=&\int_{t}^{\tau_{n}}\left(\mathcal{A}^{\hat{\Pi}}f^{s}\right)\left(\hat{X}_{r},r,\hat{\xi}_{r}\right)\mathrm{d}r+\int_{t}^{\tau_{n}}f^{s}_{x}\left(\hat{X}_{r},r,\hat{\xi}_{r}\right)\sigma\left(\hat{X}_{r},r,\hat\pi_r,\hat{\xi}_{r}\right)\mathrm{d}B_{r}\\
&+\int_{t}^{\tau_{n}}\left[f^{s}_{y}\left(\hat{X}_{r-},r,\hat{\xi}_{r-}\right)-f^{s}_{x}\left(\hat{X}_{r-},r,\hat{\xi}_{r-}\right)\right]\mathrm{d}\hat{\xi}_{r}^{c}\\
&+\sum\limits_{r\in[t,\tau_{n}]}\int_{0}^{\Delta\hat{\xi}_{r}}\left[f^{s}_{y}\left(\hat{X}_{r-}-a,r,\hat{\xi}_{r-}+a\right)-f^{s}_{x}\left(\hat{X}_{r-}-a,r,\hat{\xi}_{r-}+a\right)\right]\mathrm{d}a.
\end{aligned}
\label{fsito}\end{equation}
As $\left(\hat{X}_{r},r,\hat{\xi}_{r}\right)\in\overline{W^{\hat{\Xi}}}, \forall r\in[t,\tau^{x,t,y;\hat{\Pi},\hat{\Xi}}_{t}]$, using \eqref{fw}, we have $\left(\mathcal{A}^{\hat{\Pi}}f^{s}\right)\left(\hat{X}_{r},r,\hat{\xi}_{r}\right)=-\beta(r-s)H\left(\hat{X}_{r},r,\hat\pi_r,\hat{\xi}_{r}\right),  \forall r\in[t,\tau_n]$. Furthermore, using \eqref{martingale}, we have 
\begin{equation}
\mathbb{E}_{x,t,y}\left[\int_{t}^{\tau_{n}}f^{s}_{x}\left(\hat{X}_{r},r,\hat{\xi}_{r}\right)\sigma\left(\hat{X}_{r},r,\hat\pi_r,\hat{\xi}_{r}\right)\mathrm{d}B_{r}\right]=0.
\nonumber\end{equation}
Note that for any $r\in[t,\tau^{x,t,y;\hat{\Pi},\hat{\Xi}}_{t}]$ such that $\Delta\hat{\xi}_{r}>0$ and any $a\in[0,\Delta\hat{\xi}_{r}]$, it holds almost surely that $(\hat{X}_{r-}-a,r,\hat{\xi}_{r-}+a)\in P^{\hat{\Xi}}$. Combining with the definition of $P^{\hat{\Xi}}$, using \eqref{cxy} and \eqref{fp}, we have 
\begin{equation}\begin{aligned}
&\int_{t}^{\tau_{n}}\left[f^{s}_{y}\left(\hat{X}_{r-},r,\hat{\xi}_{r-}\right)-f^{s}_{x}\left(\hat{X}_{r-},r,\hat{\xi}_{r-}\right)\right]\mathrm{d}\hat{\xi}_{r}^{c}\\
&+\sum\limits_{r\in[t,\tau_{n}]}\int_{0}^{\Delta\hat{\xi}_{r}}\left[f^{s}_{y}\left(\hat{X}_{r-}-a,r,\hat{\xi}_{r-}+a\right)-f^{s}_{x}\left(\hat{X}_{r-}-a,r,\hat{\xi}_{r-}+a\right)\right]\mathrm{d}a\\=&
-\int_{t}^{\tau_{n}}\beta(r-s)c\left(\hat{X}_{r-},r,\hat{\xi}_{r-}\right)\mathrm{d}\hat{\xi}_{r}.
\end{aligned}\nonumber\end{equation}

Therefore, based on the above analysis, taking expectations $\mathbb{E}_{x,t,y}$ on both sides of \eqref{fsito}, letting $n\uparrow+\infty$, using \eqref{finfty1} and the dominated convergence theorem, we obtain the probabilistic interpretation \eqref{intf}.

{\bf Step 2:} We show that
\begin{equation}
\label{step2eq}
V(x,t,y)=f(x,t,y,t)=J\left(x,t,y;\hat{\Pi},\hat{\Xi}\right),\ \forall(x,t,y)\in\mathcal{Q}.
\end{equation}

Similar to {\bf Step 1}, fix any $(x,t,y)\in\mathcal{Q}$ and $n\ge t$, applying the It\^{o}--Tanaka--Meyer formula to $V$, taking expectations $\mathbb{E}_{x,t,y}$, and using \eqref{cxy} and \eqref{martingale}, we have 
\begin{equation}
\label{intv}
V(x,t,y)=\mathbb{E}_{x,t,y}\left[\int_t^{\tau_{n}} -\left(\mathcal{A}^{\hat\Pi}V\right)\left(\hat{X}_r, r, \hat{\xi}_r \right) \mathrm{d} r+V\left(\hat{X}_{\tau_{n}}, \tau_{n},\hat{\xi}_{\tau_{n}}\right)+\int_t^{\tau_{n}} c\left(\hat{X}_{r-}, r, \hat{\xi}_{r-}\right) \mathrm{d} \hat{\xi}_r \right].
\end{equation}
Again, applying the It\^{o}--Tanaka--Meyer formula to $(x,t,y)\mapsto f(x,t,y,t)$, taking expectations $\mathbb{E}_{x,t,y}$, and using \eqref{cxy}, \eqref{fp} and \eqref{martingale}, we have
\begin{equation}
\label{intft}
f(x,t,y,t)=\mathbb{E}_{x,t,y}\left[\int_t^{\tau_{n}} -\left(\mathcal{A}^{\hat\Pi}f\right)\left(\hat{X}_r, r, \hat{\xi}_r,r \right) \mathrm{d} r+ f\left(\hat{X}_{\tau_{n}}, \tau_{n},\hat{\xi}_{\tau_{n}},\tau_{n}\right)+\int_t^{\tau_{n}} c\left(\hat{X}_{r-}, r, \hat{\xi}_{r-}\right) \mathrm{d} \hat{\xi}_r \right].
\end{equation}
Since $\hat{\Pi}$ is continuous on $\overline{W^{\hat{\Xi}}}-W^{\hat{\Xi}}$  when $\hat{\Pi}$ is regarded as a function on $\overline{W^{\hat{\Xi}}}$, the combination of \eqref{v} and \eqref{fw} yields that
\begin{equation}
\label{vf}
\left(\mathcal{A}^{\hat\Pi}V\right)(x_0,t_0,y_0)=\left(\mathcal{A}^{\hat\Pi}f\right)(x_0,t_0,y_0,t_0),\ \forall (x_0,t_0,y_0)\in \overline{W^{\hat{\Xi}}}.
\end{equation}
Then, based on \eqref{intv}--\eqref{vf}, letting $n\uparrow+\infty$ and using \eqref{finfty2}, we obtain $V(x,t,y)=f(x,t,y,t)$. The second equality in \eqref{step2eq} follows directly from the probabilistic interpretation \eqref{intf}. Hence \eqref{step2eq} holds.

{\bf Step 3:} We prove that $(\hat{\Pi},\hat{\Xi})$ is an equilibrium regular-singular control law.

Fixed any $(x,t,y)\in\mathcal{Q}$. For any admissible regular-singular control pair $(\pi,\xi)\in\mathcal{D}^{x,t,y}$, we define its value function discounted to $s$ by

\begin{equation}\label{fpixi}
\begin{aligned}
f^{\pi,\xi}(x,t,y,s):=&\mathbb{E}_{x,t,y}\left[\int_t^{\tau^{x,t,y;\pi,\xi}_{t}} \beta(r-s) H\left(X^{\pi,\xi}_r, r,\pi_r, \xi_r \right) \mathrm{d} r\right. \\
& \left.+\int_t^{\tau^{x,t,y;\pi,\xi}_{t}} \beta(r-s) c\left(X^{\pi,\xi}_{r-}, r, \xi_{r-}\right) \mathrm{d}\xi_r \right],\ \forall  s\in[0,t].
\end{aligned}
\end{equation}
Clearly, we have $f^{\hat\pi,\hat\xi}(x,t,y,s)=f(x,t,y,s), \forall s\in [0,t]$. 

Fix any admissible pair of perturbations $(u,\eta)\in\mathcal{\tilde D}^{x,t,y;\hat\Pi,\hat\Xi}$, there exist constants $\tilde{h}>0$ and $\tilde{M}>0$ such that $\eta_{(t+h)-}-\eta_{t}\le \tilde{M}h, \text{a.s.}, \forall h\in(0,\tilde{h})$, and the perturbed regular-singular control pair $\left(\pi^{h},\xi^{h}\right)$, $\forall h\in(0,\tilde{h})$ defined in \eqref{pihxih} and the corresponding state process $X^{\pi^{h},\xi^{h}}$ in \eqref{xpihxih} satisfy integrability condition (b) in Definition \ref{admissible2}. Thus, we can define $f^{\pi^{h},\xi^{h}}(x,t,y,s), \forall h\in(0,\tilde{h}), s\in[0,t]$ just as \eqref{fpixi}. 

Moreover, by \eqref{pihxih}, \eqref{xpihxih} and \eqref{fT}, we can define 
\begin{equation}\begin{aligned}
&f^{\pi^{h},\xi^{h}}\left(X^{\pi^{h},\xi^{h}}_{\tau^{x,t,y;u,\eta}_{t}\wedge((t+h)-)},\tau^{x,t,y;u,\eta}_{t}\wedge(t+h),\xi^{h}_{\tau^{x,t,y;u,\eta}_{t}\wedge((t+h)-)},s\right)\\=&f^{\pi^{h},\xi^{h}}\left(X^{u,\eta}_{\tau^{x,t,y;u,\eta}_{t}\wedge((t+h)-)},\tau^{x,t,y;u,\eta}_{t}\wedge(t+h),\eta_{\tau^{x,t,y;u,\eta}_{t}\wedge((t+h)-)},s\right)\\=&f\left(X^{u,\eta}_{\tau^{x,t,y;u,\eta}_{t}\wedge((t+h)-)},\tau^{x,t,y;u,\eta}_{t}\wedge(t+h),\eta_{\tau^{x,t,y;u,\eta}_{t}\wedge((t+h)-)},s\right),\ \forall s\in[0,t], h\in(0,\tilde{h}).
\end{aligned}\nonumber\end{equation}
%Besides, we know that $\tau^{x,t,y;\pi^{h},\xi^{h}}_{t}=\tau^{x,t,y;u,\eta}_{t}$, when $\tau^{x,t,y;u,\eta}_{t}<t+h$, and $1_{\left\{\tau^{x,t,y;\pi^{h},\xi^{h}}_{t}\ge t+h\right\}}=1_{\left\{\tau^{x,t,y;u,\eta}_{t}\ge t+h\right\}}, \forall h\in(0,\tilde{h})$. Therefore, $\tau^{x,t,y;\pi^{h},\xi^{h}}\wedge(t+h)=\tau^{x,t,y;u,\eta}\wedge(t+h), \forall h\in(0,\tilde{h})$.

Hereafter, we assume $h$ is arbitrary in $(0,\tilde{h}\wedge h_0)$, where $h_0>0$ is the constant in condition (6). Observing that $\pi^{h}_{r}=u_{r}, \xi^{h}_{r}=\eta_{r}, X^{\pi^{h},\xi^{h}}_r=X^{u,\eta}_r, \forall r\in[t,\tau^{x,t,y;\pi^{h},\xi^{h}}_{t}\wedge((t+h)-)]$, we have 
\begin{equation}
\begin{aligned}
J\left(x,t,y;\pi^{h},\xi^{h}\right)=&\mathbb{E}_{x,t,y}\left[f^{\pi^{h},\xi^{h}}\left(X^{u,\eta}_{\tau^{x,t,y;u,\eta}_{t}\wedge((t+h)-)},\tau^{x,t,y;\pi^{h},\xi^{h}}_{t}\wedge(t+h),\eta_{\tau^{x,t,y;u,\eta}_{t}\wedge((t+h)-)},t\right)\right.\\
  &\left.+\int_{t}^{\tau^{x,t,y;u,\eta}_{t}\wedge(t+h)}\beta(r-t) H\left(X^{u,\eta}_r, r,u_r, \eta_r \right) \mathrm{d} r \right.\\
  &\left.+\int_{t}^{\tau^{x,t,y;u,\eta}_{t}\wedge(t+h)}\beta(r-t) c\left(X^{u,\eta}_{r-}, r, \eta_{r-}\right) \mathrm{d}\eta^{c}_r\right.\\
  &\left.+\sum\limits_{r\in[t,\tau^{x,t,y;u,\eta}_{t}\wedge((t+h)-)]}\beta(r-t) c\left(X^{u,\eta}_{r-}, r, \eta_{r-}\right) \Delta\eta_r\right].
\end{aligned}\nonumber\end{equation}
Then, using \eqref{step2eq} and the fact that $f^{\pi^{h},\xi^{h}}\left(X^{u,\eta}_{\tau^{x,t,y;u,\eta}_{t}\wedge((t+h)-)},\tau^{x,t,y;u,\eta}_{t}\wedge(t+h),\eta_{\tau^{x,t,y;u,\eta}_{t}\wedge((t+h)-)},t\right)\\=f\left(X^{u,\eta}_{\tau^{x,t,y;u,\eta}_{t}\wedge((t+h)-)},\tau^{x,t,y;u,\eta}_{t}\wedge(t+h),\eta_{\tau^{x,t,y;u,\eta}_{t}\wedge((t+h)-)},t\right)$, we obtain
\begin{equation}
\begin{aligned}
&J\left(x,t,y;\pi^{h},\xi^{h}\right)-J\left(x,t,y;\hat\pi,\hat\xi\right)\\
=&\mathbb{E}_{x,t,y}\left[f\left(X^{u,\eta}_{\tau^{x,t,y;u,\eta}_{t}\wedge((t+h)-)},\tau^{x,t,y;u,\eta}_{t}\wedge(t+h),\eta_{\tau^{x,t,y;u,\eta}_{t}\wedge((t+h)-)},t\right)-V(x,t,y)\right.\\
  &\left.+\int_{t}^{\tau^{x,t,y;u,\eta}_{t}\wedge(t+h)}\beta(r-t) H\left(X^{u,\eta}_r, r,u_r, \eta_r \right) \mathrm{d} r \right.\\
  &\left.+\int_{t}^{\tau^{x,t,y;u,\eta}_{t}\wedge(t+h)}\beta(r-t) c\left(X^{u,\eta}_{r-}, r, \eta_{r-}\right) \mathrm{d}\eta^{c}_r\right.\\
  &\left.+\sum\limits_{r\in[t,\tau^{x,t,y;u,\eta}_{t}\wedge((t+h)-)]}\beta(r-t) c\left(X^{u,\eta}_{r-}, r, \eta_{r-}\right) \Delta\eta_r\right].
\end{aligned}\nonumber\end{equation}

Similar to {\bf Step 1}, applying the It\^{o}--Tanaka--Meyer formula to $V$, taking expectations $\mathbb{E}_{x,t,y}$, and using \eqref{martingale2} and \eqref{step2eq}, we deduce that
\begin{equation}
\begin{aligned}
&\mathbb{E}_{x,t,y}\left[f\left(X^{u,\eta}_{\tau^{x,t,y;u,\eta}_{t}\wedge((t+h)-)},\tau^{x,t,y;u,\eta}_{t}\wedge(t+h),\eta_{\tau^{x,t,y;u,\eta}_{t}\wedge((t+h)-)},t\right)-V(x,t,y)\right]\\
=&\mathbb{E}_{x,t,y}\left[f\left(X^{u,\eta}_{\tau^{x,t,y;u,\eta}_{t}\wedge((t+h)-)},\tau^{x,t,y;u,\eta}_{t}\wedge(t+h),\eta_{\tau^{x,t,y;u,\eta}_{t}\wedge((t+h)-)},t\right)\right.\\ &\left.-f\left(X^{u,\eta}_{\tau^{x,t,y;u,\eta}_{t}\wedge((t+h)-)},\tau^{x,t,y;u,\eta}_{t}\wedge(t+h),\eta_{\tau^{x,t,y;u,\eta}_{t}\wedge((t+h)-)},\tau^{x,t,y;u,\eta}_{t}\wedge(t+h)\right)\right.\\ &\left.+ V\left(X^{u,\eta}_{\tau^{x,t,y;u,\eta}_{t}\wedge((t+h)-)},\tau^{x,t,y;u,\eta}_{t}\wedge(t+h),\eta_{\tau^{x,t,y;u,\eta}_{t}\wedge((t+h)-)}\right)-V(x,t,y)\right]\\
=&\mathbb{E}_{x,t,y}\left[\int_{t}^{\tau^{x,t,y;u,\eta}_{t}\wedge(t+h)}\left[\left(\mathcal{A}^{u}V\right)\left(X^{u,\eta}_{r},r,\eta_{r}\right)\!-\!f_s\left(X^{u,\eta}_{\tau^{x,t,y;u,\eta}_{t}\wedge((t+h)-)},\tau^{x,t,y;u,\eta}_{t}\!\wedge\!(t\!+\!h),\eta_{\tau^{x,t,y;u,\eta}_{t}\wedge((t+h)-)},r\right)\right]\mathrm{d}r\right.\\ &\left.+\int_{t}^{\tau^{x,t,y;u,\eta}_{t}\wedge(t+h)}\left[V_{y}\left(X^{u,\eta}_{r-},r,\eta_{r-}\right)-V_{x}\left(X^{u,\eta}_{r-},r,\eta_{r-}\right)\right]\mathrm{d}\eta_{r}^{c}\right.\\
&\left.+\sum\limits_{r\in[t,\tau^{x,t,y;u,\eta}_{t}\wedge((t+h)-)]}\int_{0}^{\Delta\eta_{r}}\left[V_{y}\left(X^{u,\eta}_{r-}-a,r,\eta_{r-}+a\right)-V_{x}\left(X^{u,\eta}_{r-}-a,r,\eta_{r-}+a\right)\right]\mathrm{d}a\right],
\end{aligned}\nonumber\end{equation}
where $u$ in $\left(\mathcal{A}^{u}V\right)$ refers to the regular perturbation process $u$, i.e.,
\begin{equation}\left(\mathcal{A}^{u}\varphi\right)(x,t,y):=\varphi_{t}(x,t,y)+\mu(x,t,u_t,y)\varphi_{x}(x,t,y)+\frac{1}{2}\sigma^{2}(x,t,u_t,y)\varphi_{xx}(x,t,y),\ \forall\varphi:\mathcal{Q}\rightarrow\mathbb{R}.\nonumber\end{equation}

Therefore, combining with \eqref{cxy}, the definition of $\hat\Pi$, and \eqref{v}, we have
\begin{equation}
\begin{aligned}
&J\left(x,t,y;\pi^{h},\xi^{h}\right)-J\left(x,t,y;\hat\pi,\hat\xi\right)\\=&\mathbb{E}_{x,t,y}\left[\int_{t}^{\tau^{x,t,y;u,\eta}_{t}\wedge(t+h)}\left[\left(\mathcal{A}^{u}V\right)\left(X^{u,\eta}_{r},r,\eta_{r}\right)+\beta(r-t) H\left(X^{u,\eta}_r, r,u_r, \eta_r \right) \right]\mathrm{d} r\right.\\
  &\left.-\int_{t}^{\tau^{x,t,y;u,\eta}_{t}\wedge(t+h)}f_s\left(X^{u,\eta}_{\tau^{x,t,y;u,\eta}_{t}\wedge((t+h)-)},\tau^{x,t,y;u,\eta}_{t}\wedge(t+h),\eta_{\tau^{x,t,y;u,\eta}_{t}\wedge((t+h)-)},r\right)\mathrm{d}r \right.\\
  &\left.+\int_{t}^{\tau^{x,t,y;u,\eta}_{t}\wedge(t+h)}\left[\beta(r-t) c\left(X^{u,\eta}_{r-}, r, \eta_{r-}\right)+\left(V_{y}-V_{x}\right)\left(X^{u,\eta}_{r-},r,\eta_{r-}\right)\right] \mathrm{d}\eta^{c}_r\right.\\
  &\left.+\sum\limits_{r\in[t,\tau^{x,t,y;u,\eta}_{t}\wedge((t+h)-)]}\int_{0}^{\Delta\eta_{r}}\left[\beta(r-t) c\left(X^{u,\eta}_{r-}-a, r, \eta_{r-}+a\right)+\left(V_{y}-V_{x}\right)\left(X^{u,\eta}_{r-}-a,r,\eta_{r-}+a\right)\right]\mathrm{d}a\right]\\
  \ge&\mathbb{E}_{x,t,y}\left[\int_{t}^{\tau^{x,t,y;u,\eta}_{t}\wedge(t+h)}\left[f_s\left(X^{u,\eta}_{r},r,\eta_{r},r\right)+(\beta(r-t)-1) H\left(X^{u,\eta}_r, r,u_r, \eta_r \right) \right]\mathrm{d} r\right.\\
  &\left.-\int_{t}^{\tau^{x,t,y;u,\eta}_{t}\wedge(t+h)}f_s\left(X^{u,\eta}_{\tau^{x,t,y;u,\eta}_{t}\wedge((t+h)-)},\tau^{x,t,y;u,\eta}_{t}\wedge(t+h),\eta_{\tau^{x,t,y;u,\eta}_{t}\wedge((t+h)-)},r\right)\mathrm{d}r \right.\\
  &\left.+\int_{t}^{\tau^{x,t,y;u,\eta}_{t}\wedge(t+h)}(\beta(r-t)-1) c\left(X^{u,\eta}_{r-}, r, \eta_{r-}\right)\mathrm{d}\eta^{c}_r\right.\\
  &\left.+\sum\limits_{r\in[t,\tau^{x,t,y;u,\eta}_{t}\wedge((t+h)-)]}\int_{0}^{\Delta\eta_{r}}(\beta(r-t)-1) c\left(X^{u,\eta}_{r-}-a, r, \eta_{r-}+a\right)\mathrm{d}a\right].
\end{aligned}\label{est0}\end{equation}

Using \eqref{integrability} and the dominated convergence theorem, since the regular control space $U\subset\mathbb{R}$ is a compact set and $u_r\in U$, we have 
\begin{equation}\begin{aligned}
&\lim\limits_{h\downarrow 0}\mathbb{E}_{x,t,y}\left[\frac{1}{h}\left(\int_{t}^{\tau^{x,t,y;u,\eta}_{t}\wedge(t+h)}f_s\left(X^{u,\eta}_{r},r,\eta_{r},r\right)+(\beta(r-t)-1) H\left(X^{u,\eta}_r, r,u_r, \eta_r \right) \right.\right.\\
  &\left.\left.-f_s\left(X^{u,\eta}_{\tau^{x,t,y;u,\eta}_{t}\wedge((t+h)-)},\tau^{x,t,y;u,\eta}_{t}\wedge(t+h),\eta_{\tau^{x,t,y;u,\eta}_{t}\wedge((t+h)-)},r\right)\mathrm{d}r\right)\right]\\
=&\mathbb{E}_{x,t,y}\left[\left[f_s\left(X^{u,\eta}_{t},t,\eta_{t},t\right)+ (\beta(t-t)-1)H\left(X^{u,\eta}_t, t,u_t, \eta_t \right)-f_s\left(X^{u,\eta}_{t},t,\eta_{t},t\right)\right]1_{\left\{\tau^{x,t,y;u,\eta}_{t}>t\right\}} \right]=0.
\end{aligned}\label{est1}\end{equation}
Since $\eta_{(t+h)-}-\eta_{t}\le \tilde{M}h, \text{a.s.}$, we know that $\int_{t}^{t+h}\mathrm{d}\eta^{c}_{r}\le \tilde{M}h, \text{a.s.}$ and $\sum\limits_{r\in(t,t+h)}\Delta\eta_{r}\le \tilde{M}h, \text{a.s.}$. Then, using \eqref{cxy}, \eqref{integrability2} and the Fatou's lemma, we have
\begin{equation}\begin{aligned}
&\liminf\limits_{h\downarrow 0}\mathbb{E}_{x,t,y}\left[\frac{1}{h}\int_{t}^{\tau^{x,t,y;u,\eta}_{t}\wedge(t+h)}\left(\beta(r-t)-1\right) c\left(X^{u,\eta}_{r-}, r, \eta_{r-}\right) \mathrm{d}\eta^{c}_r\right]\\
\geq&\mathbb{E}_{x,t,y}\left[\liminf\limits_{h\downarrow 0}\frac{1}{h}\int_{t}^{\tau^{x,t,y;u,\eta}_{t}\wedge(t+h)}\left(\beta(r-t)-1\right) c\left(X^{u,\eta}_{r}, r, \eta_{r}\right)\mathrm{d}\eta^{c}_r \right]\\
=&\mathbb{E}_{x,t,y}\left[ \liminf\limits_{h\downarrow 0}\frac{1}{h}\int_{t}^{\tau^{x,t,y;u,\eta}_{t}\wedge(t+h)}\left(\beta(t-t)-1\right) c(X^{u,\eta}_t, t, \eta_t)\mathrm{d}\eta^{c}_{r}\right]= 0,
\end{aligned}\label{est2}\end{equation}
and
\begin{equation}\begin{aligned}
&\liminf\limits_{h\downarrow 0}\mathbb{E}_{x,t,y}\left[\frac{1}{h}\sum\limits_{r\in[t,\tau^{x,t,y;u,\eta}_{t}\wedge((t+h)-)]}\int_{0}^{\Delta\eta_{r}}\left(\beta(r-t)-1\right) c\left(X^{u,\eta}_{r-}-a, r, \eta_{r-}+a\right)\mathrm{d}a\right]\\
=&\liminf\limits_{h\downarrow 0}\mathbb{E}_{x,t,y}\left[\frac{1}{h}\sum\limits_{r\in(t,\tau^{x,t,y;u,\eta}_{t}\wedge((t+h)-)]}\left(\beta(r-t)-1\right) c\left(X^{u,\eta}_{r-}, r, \eta_{r-}\right)\Delta\eta_{r}\right]\\
\geq&\mathbb{E}_{x,t,y}\left[\liminf\limits_{h\downarrow 0}\frac{1}{h}\sum\limits_{r\in(t,\tau^{x,t,y;u,\eta}_{t}\wedge((t+h)-)]}\left(\beta(r-t)-1\right) c\left(X^{u,\eta}_{r}, r, \eta_{r}\right)\Delta\eta_{r} \right]\\
=&\mathbb{E}_{x,t,y}\left[\liminf\limits_{h\downarrow 0}\frac{1}{h}\sum\limits_{r\in(t,\tau^{x,t,y;u,\eta}_{t}\wedge((t+h)-)]}\left(\beta(t-t)-1\right)c(X^{u,\eta}_t, t, \eta_t)\Delta\eta_{r}\right]= 0.
\end{aligned}\label{est3}\end{equation}
\eqref{est0}--\eqref{est3} yield that the equilibrium condition \eqref{equicond} holds. Hence, $(\hat{\Pi},\hat{\Xi})$ is an equilibrium  regular-singular control law.
\end{proof}

\end{document}
